\scanpage{25}
Soit $A^+$ son cône positif, \deleted{engendré} adhérence
de l'ens des combinaisons linéaires à
coefficients $\geq0$ des $u^\varphi u$ ($u\in A$ ; on
peut $=$ supposer $u\in A^*$, groupe des \uline{unités} de $A$)

Il correspond dans $\Phi^s$ au\deleted{\unsure{x}} cône fermé $\Phi^+$
engendré par les $\varphi^u$ $[\varphi^u(x,y)=\varphi(ux,uy)]$,
i.e. le plus petit cône fermé dans $\Phi^s$
contenant $\varphi$ et stable par les automorphismes
de $M$. \deleted{Tout élément de} De la théorie des
Algèbres de Riemann résulte que tt
élément $a$ \uline{intérieur} à $A^+$ est une
unité, donc tt élément intérieur à $\Phi^+$
est une forme non dégénérée ; \deleted{\unsure{Donc on}} \inserted{de plus,}
pour un tel $a$, \deleted{\unsure{l'involut|l'involution}} (correspondant à
\deleted{telle que \unsure{engendre} \unsure{le|une des} $=$ cône}
$\varphi\in\Phi^+$ intérieur à $\Phi^+$) l'involution
$u\mapsto u^\psi=a^{-1}u^\varphi a$ \deleted{satisfait à} étant déduite
\[
\deleted{\unsure{\mathrm{Tr}}\,u^\psi u=\unsure{\mathrm{Tr}}\,a^{-1}u^\varphi au=\sum_i b_i^{\deleted{\unsure{*}}-1}(b_i^\varphi)^{-1}u^\varphi b_i^\varphi b_i u}
\]
de la précédente par un autom (intérieur)
qui ne bouge pas la trace, elle \deleted{\unsure{peut}} possède
% archive sheet number: 24
